%!TeX root=csp-audit.tex
\section{Steps performed silently}\label{sec:expansion}

These are not defects. They are the places where the proof paper must write out
an argument that \cite{ZhukSimplified} performs in a clause, and they are the
bulk of the work: about two dozen of them against fifteen register entries.
They are collected here so that a reader comparing the two documents can see
where the extra text comes from and satisfy themselves that it is expansion
rather than divergence.

\subsection{Coverage of \S5}\label{sec:coverage}

\S5 of \cite{ZhukSimplified} is $3\,144$ lines, $44$ statements with proofs,
$152$ proof-dependency edges. All seven subsections have now been read.

\begin{center}\small
\begin{tabular}{@{}llll@{}}
\hline
Lines & Subsection & Proofs & Defects found \\
\hline
48--70 & Additional definitions & 0 & --- \\
71--346 & Subuniverses of types $\TBA$, $\TC$, $\TS$ & 6 & none \\
347--1021 & Intersection property & 6 & none; $\mathrm F$ localised here \\
1022--1735 & Properties of PC or linear congruences & 9 & $\mathrm D11$, $\mathrm D12$, $\mathrm D13$, $\mathrm D14$ \\
1736--1927 & Types interaction & 2 & none; one missing mirror case \\
1928--2296 & Factorisation of strong subalgebras & 5 & $\mathrm D15$ \\
2297--3144 & Proof of the remaining statements & 16 & $\mathrm D16$, $\mathrm D17$ \\
\hline
\end{tabular}
\end{center}

The distribution is not an accident. \S5.4 is the only part of \S5 that
manipulates bridges as objects in their own right rather than consuming them,
and every one of its four defects is about which auxiliary property survives a
construction --- symmetrisation, restriction to a block, restriction to a
strong subuniverse.

\subsection{The four hedges}\label{sec:hedges}

A mechanical sweep for hedging phrases over the thirteen printed pages of \S5
returns four, an unusually low density: the same regex over other expositions
in this subject fires many times per page. \textbf{All four are honest.}

\begin{center}\footnotesize
\begin{tabular}{@{}lp{0.24\textwidth}p{0.50\textwidth}@{}}
\hline
Line & In & Disposition \\
\hline
\texttt{SS:146} & \texttt{LEMBACenterSImplyFactor} &
  ``For $T=\TBA$ it is straightforward'': it is one line --- if
  $t(B,A)\subseteq B$ and $t(A,B)\subseteq B$ then the same $t$ works on
  images. Only $T=\TC$ is genuinely imported. \\
\texttt{SS:177} & \texttt{LEMBACenter\brk SOnPower\brk Implies} &
  ``For $T=\TS$ just repeat the same proof \emph{word to word}'': no repetition
  is needed --- see below. \\
\texttt{SS:1560} & before \texttt{LEMNoBridge\brk Between\brk DifferentTypes} &
  ``This case can be considered in the same way as Case 2'': it is genuinely
  the mirror. $(a,b,a,b)\in\xi$ places $b$ in the block and
  $\cover\sigma\subseteq\operatorname{RightLinked}$ places $a$ and $c$;
  writing it out is mechanical. \\
\texttt{SS:2428} & \texttt{LEMFactorByDelta} region &
  ``The inclusion $\subseteq$ is obvious'': it is. \\
\hline
\end{tabular}
\end{center}

\begin{entry}[The source undersells its own citation]\label{ent:6.24}
\texttt{LEMBACenterSOnPowerImplies} says $B\subt{T}\mathbf A^{n}$ with
$T\in\{\TBA,\TC,\TS\}$ implies some $C\subt{T}\mathbf A$, and proves it by
``For $T\in\{\TBA,\TC\}$ see \cite[Lem.~6.24]{ZhukStrong}. For $T=\TS$ just
repeat the same proof word to word replacing $\TBA$ by $\TS$.'' The cited lemma
reads: \emph{``Suppose $R$ is a nontrivial strong subuniverse of $A_{1}\times
\dots\times A_{n}$ of type $T\neq\TPC$. Then there exists $i$ such that $A_{i}$
has a nontrivial subuniverse of type $T$.''} It is already general in $T$ and
already for a product of distinct algebras --- the power is a special case ---
and its proof is \emph{type-uniform}: it takes $\proj_{1}(R)$ if that is proper
and otherwise a fibre, and neither construction mentions the type, which enters
only through the facts that projections and fibres preserve $\TBA$ and preserve
$\TC$ respectively. So a single run of the proof produces one witness carrying
\emph{every} strong type $R$ has at once; if $R$ is simultaneously binary
absorbing and central --- which is what $\TS$ means --- the witness is too. The
hedge disappears rather than being discharged three times. Proof paper:
\pkey{lem:power-to-base}, \pkey{haz:power-to-base}.
\end{entry}

\subsection{\S5.3, the intersection property}\label{sec:exp-53}

\texttt{LEMIntersectionPCLinearIsGood} (\texttt{SS:623}) is a 292-line proof
with 20 distinct citations, by a wide margin the largest and most connected in
\S5. It was read in full and is \textbf{structurally sound}. Every appeal to
the inductive hypothesis lands at strictly smaller $k+\ell$: part (s) uses (s)
and (d) at $k+(\ell-1)$; subsubcase 2B2 uses (s) at $(n-1)+k$; 2B3 uses (d) at
$(n-1)+k$; and $n\le\ell$ throughout. That is worth confirming, since a
commented-out block shows an earlier draft inducted lexicographically on
$(|A/(\sigma\cap\omega)|,\,k+\ell)$ and the live proof dropped the first
component.

Seven items for the rendering. The proof paper now carries the expansion in
full; \pkey{haz:intersection-draft} keeps the four of these that a reader of
the finished proof still needs. Every consumer of the lemma was re-read against
the written-out statement afterwards --- \pkey{cor:intersection-good},
\pkey{lem:parallelogram-D}, \pkey{lem:intersect-all}(t),
\pkey{thm:stable-intersection} in its no-$\TS$, mixed and both-$\TD$ branches,
and \pkey{lem:two-stable}(2) --- and none of them breaks. Two points are worth
recording. No call site consumes minimality of the chains, which item~5 below
drops. And \pkey{lem:intersect-all}(t) is proved \emph{from} the lemma, in both
parts, so it may not be used inside it; the proof in the proof paper does not.

\begin{enumerate}\tightlist
\item \emph{Part (s): five properness splits, not one.} Throughout part (s) the
  text asserts \emph{proper} containments where
  \texttt{LEMBACenterImplyIntersection} and \texttt{LEMBACenterSImplyFactor}
  give only $\subseteq$. None is a defect in the lemma, because every equality
  branch gives the conclusion outright, but each needs writing. Branch
  $\mathcal T_{\ell}=\TD$: if
  $(G\cap C_{\ell-1})/\omega_{\ell}=C_{\ell-1}/\omega_{\ell}$ then every
  $\omega_{\ell}$-block meeting $C_{\ell-1}$ meets $G\cap C_{\ell-1}$, in
  particular the block $E$ with $C_{\ell}=C_{\ell-1}\cap E$, so
  $G\cap C_{\ell}\neq\varnothing$. Branches
  $\mathcal T_{\ell}\in\{\TBA,\TC\}$ and $\mathcal T_{\ell}=\TS$ each carry two
  more, one per containment into $B\cap C_{\ell-1}$; there the equality branch
  makes one of the two sets contain the other and the conclusion is immediate.
\item \emph{What Lemma 6.25 is doing in part (s).} Part (s) appeals to
  \texttt{LEMBACenterS\brk Possible\brk Intersections} to conclude
  nonemptiness, from a lemma whose conclusion is ``$T_{1}=T_{2}\in\{\TBA,\TC\}$'', which by itself
  permits two disjoint binary absorbing subuniverses. The step is valid only
  because $\subt{\TBA,\TC}$ is a \emph{conjunction} (\texttt{main.tex:1567}):
  one of the two sets carries both types, so the type opposite to the other
  set's is available and equality of types fails. Proof paper:
  \pkey{lem:ba-central-meets}, \pkey{haz:ba-central-meets}.
\item \emph{Properness in 2A1 and 2B1} comes from the minimal choice of $m$
  (resp.\ $n$); nonemptiness from the constant tuples
  $(b/\delta,\dots,b/\delta)$ for $b\in B$, which lie in $S_{m,0}$ since
  $B\subseteq B_{m}$ and $(b,b)\in\sigma\cap\omega$.
\item \emph{Subsubcase 2B3} writes the first alternative as
  ``$=C_{n}/\omega_{n}$'', which is stronger than the ``size 1'' the inductive
  hypothesis supplies. Correct, because $B\cap C\neq\varnothing$ and
  $C\subseteq C_{n}$ force $B_{k}\cap C_{n}\neq\varnothing$, so the single
  block is the one defining $C_{n}$; then $B_{k}\cap C_{n-1}\subseteq C_{n}$
  and $S_{k,n-1}$, $S_{k,n}$ cut $(B/\delta)^{|A|}$ identically, contradicting
  minimality of $n$.
\item \emph{Subsubcase 2A3 reads one alternative off a two-alternative lemma.}
  ``\texttt{LEMSelf\brk IntersectionPC} implies that
  $((B\circ\delta)\cap B_{m-1})/\sigma_{m}=B_{m-1}/\sigma_{m}$'' (\texttt{SS:828})
  is not an implication: that lemma offers $B_{m-1}/\sigma_{m}$ \emph{or}
  $B_{m}/\sigma_{m}$. The second is excluded, but by an argument the text does
  not give: $B_{m}=B_{m-1}\cap E$ lies inside one $\sigma_{m}$-class, so the
  second alternative makes the last coordinate of $R'$ constant, and with
  $\proj_{1,\dots,|A|}(R')=(B/\delta)^{|A|}$ that yields
  $(B/\delta)^{|A|}\times B_{m}/\sigma_{m}\subseteq R'$ --- which the minimal
  choice of $m$ forbids. Sound, but a real step.
\item \emph{The full-fibre step} at \texttt{SS:545, 837, 921} is register entry
  $\mathrm F$ (\S\ref{sec:F}).
\item \emph{Is minimality of the chains $k,\ell$ used? No.} \textbf{Settled,
  and the hypothesis is dropped.} Writing the proof out decides it in the
  negative and shows it could not have gone the other way: the inductive
  appeals in part (s), 2B2 and 2B3 are to \emph{truncations} of the chain for
  $C$, and a truncation of a minimal chain need not be minimal, so a step
  consuming minimality would break the induction rather than use it.
\item \emph{The corollary} writes $\delta=f^{-1}(\sigma)$ with $f$ undefined;
  read $f_{1}$. Its conclusion admits ``empty'' where the lemma requires
  $B\cap C\neq\varnothing$, so the empty case needs its own line at every call
  site. Proof paper: \pkey{haz:cor-intersection}.
\end{enumerate}

\texttt{LEMSelfIntersectionPC} (\texttt{SS:419}), which subsubcase 2A3
consumes, carries a proof of the same shape and three more silent steps. (i)
Subcase 1B displays ``$(E\cap B_{k-1})/\sigma_{k}=B_{k-1}/\sigma_{k-1}$''
(\texttt{SS:481}); the right-hand side is an index slip for
$B_{k-1}/\sigma_{k}$ --- $\sigma_{k-1}$ has no role in that subcase --- which
is the typo the external review flagged. (ii) The same branch concludes
$((B_{k}\circ\delta)\cap B_{k-1})/\sigma_{k}=B_{k-1}/\sigma_{k}$ from that
identity, which needs $E$ to \emph{meet} $B_{k}$; it does, because the
$\sigma_{k}$-classes the witnesses exhaust include the class of $B_{k}$, but
nothing in the text says so. (iii) Subcase 2C proves $R$ full --- a statement
about $B_{k-1}$ --- and then asserts the conclusion, which is about $B_{k}$.
The transfer works because $B_{k}=B_{k-1}\cap E_{k}$ with $E_{k}$ a
$\sigma_{k}$-class, so a $\sigma_{k}$-move inside $B_{k-1}$ cannot leave
$B_{k}$. Proof paper: \pkey{lem:self-intersection-pc},
\pkey{haz:self-intersection}.

Three smaller lemmas in the same subsection, never opened before the final
pass, are all sound. \texttt{LEMTotallySymmetricWithoutBACenter}
(\texttt{SS:357}) compresses one step: ``by the conditions of the lemma we have
$(a,\dots,a,c),(a,\dots,a,a)\in R$'' does not follow from one application of
the shift, which deletes the \emph{last} entry and so destroys $c$. What is
really happening is that symmetry plus the shift lets one pass from a tuple
with entry multiset $M$ to any $M+\{u\}-\{v\}$ with $u,v\in M$; $k-2$ such
steps delete the $b_{i}$ and duplicate $a$. This deserves its own one-line
lemma, because the same manoeuvre is needed again in
\texttt{LEMTotallySymmetricRelationForIrreducible} --- where Case~2's ``it
remains to apply the previous lemma to $R\cap(B/\sigma)^{n}$'' needs
$\proj_{1,2}(R\cap(B/\sigma)^{n})=(B/\sigma)^{2}$, not merely
$\proj_{1,2}(R)\supseteq(B/\sigma)^{2}$. Proof paper:
\pkey{lem:multiset-shift}, \pkey{haz:shift}.

\subsection{\S5.4, the bridge lemmas}\label{sec:exp-54}

Beyond $\mathrm D11$--$\mathrm D14$, \texttt{LEMBlockOfGoodBridgeDoesNotHaveBAC}
is sound in both cases, and uses five steps without stating them.

\begin{enumerate}\tightlist
\item $\bcol\delta$ is \emph{reflexive}, because
  $0_{A}\subseteq\proj_{1,2}(\delta)\subseteq\bcol\delta$, the first inclusion
  being clause (c). Everything reflexive downstream comes from this.
\item $\proj_{1,2}(\sigma\cap E^{4})=\omega\cap E^{2}$ for a subuniverse $E$,
  which is what makes Case~1's choice of $E$ the right one. It needs
  $\sigma(x_{1},x_{2},x_{1},x_{2})$ for $(x_{1},x_{2})\in\omega$ ---
  pair-reflexivity again, from the corrected symmetrisation.
\item \emph{The $E$ of Case~1 exists.} Maximality of $E$ is never used, so it
  suffices to take the block containing $a$ and $b$ of the transitive closure
  of a compatible reflexive symmetric relation on $D$ whose components are
  those of the undirected $\omega$-graph. Note that
  $(\omega\cup\omega^{-1})\cap D^{2}$, which suggests itself, is \emph{not} a
  subalgebra --- a union of two subuniverses need not be closed; use
  $(\omega\circ\omega^{-1})\cap D^{2}$, which is compatible, reflexive and
  symmetric, and has the same components because $\omega$ is reflexive.
\item The ``without loss of generality'' at \texttt{SS:1170} needs
  $\bcol\sigma$ symmetric, which holds for $\bcol\delta\circ\bcol\delta^{-1}$.
\item $\{a\}\circ\omega$ is a subuniverse, because $\{a\}$ is one by
  idempotence, which is what makes it equal to $A$ in Case~2.
\end{enumerate}

The appeal to the relational stable-intersection corollary at \texttt{SS:1214}
is \textbf{correct but not in the corollary's shape}. The three displayed boxes
$\{a\}\times B\times C\times B$, $\{a\}\times C\times C\times C$,
$\{a\}\times C\times B\times B$ are the tight boxes of three witnesses, whereas
the corollary wants $A$ in the freed coordinate. Widening each to
$\{a\}\times A\times C\times B$, $\{a\}\times C\times C\times A$,
$\{a\}\times C\times A\times B$ gives its hypothesis with $n=3$ and all
ambients $A$, legitimate because $\lll$ is reflexive; and since all three types
are $\TC$ and $n=3$, every alternative fails. The three witnesses are
$(a,b,a,b)$, $(a,a,a,a)$, $(a,a,b,b)$. Two smaller things: $C'\subte{\TC}C$ at
\texttt{SS:1226} comes out of pp-propagation as $\subte{\TC}A$ and needs
intersection to land in $C$; and the appeal there to $\proj_{1}(\xi)=A$ should
be to $\{a\}\circ\omega=A$.

\texttt{LEMNontrivialReflexiveBridgeImplies} is otherwise sound: the uniqueness
of $\cover\sigma$, the upgrade from ``every block of
$\operatorname{LeftLinked}(\cover\sigma)$ of size $>1$ satisfies
$B^{2}\subseteq\cover\sigma$'' to $\cover\sigma=
\operatorname{LeftLinked}(\cover\sigma)$, and the prime argument all check out.
Note that the contradiction is with the \emph{minimality of $\cover\sigma$},
not with irreducibility directly; the source says ``contradicts the
irreducibility of $\sigma$''. The distinction matters where $\cover\sigma$ is
data attached to a proof of irreducibility. Proof paper:
\pkey{haz:minimality-not-irred}.

\texttt{LEMPCCongruencePropertyInductiveStep} is sound outside $\mathrm D14$,
and was checked line by line. The $|A|=1$ worry does not arise: clause (c)
forces $\proj_{1,2}(\delta)\supsetneq0_{A}$, hence $|A|\ge2$. Case~2's two
applications of the ternary absorbing operation are correct and were recomputed
coordinatewise; the step needs \emph{ternary} absorption where $T=\TC$ supplies
only centrality, so ``central implies ternary absorbing''
(\texttt{main:1350}) is an unstated citation. Case~3 needs a word for the
equality case, since the absorbing-equality lemma hypothesises a proper
containment while the conclusion is immediate when it is an equality.

\texttt{LEMPCBridgesAreTrivial} is sound, with two steps to write out, both in
Case~1.

\emph{The counting step} (\texttt{SS:1478}): ``since $\proj_{1,2}(\delta)=
\proj_{3,4}(\delta)$, for any $(a,b,c,d)\in\delta$ the elements $c/\sigma$ and
$d/\sigma$ are also uniquely determined by $a/\sigma$ and $b/\sigma$''. What
the case hypothesis gives is uniqueness in the \emph{other} direction --- it is
a statement about $\delta\ast\delta^{-1}$ and the claim is about
$\delta^{-1}\ast\delta$. The bridge between them is pigeonhole, and the cited
equality of projections is what powers it: on $\cover\sigma/\sigma$, which is
finite, $\delta$ induces a bipartite relation whose left neighbourhoods are
nonempty, pairwise disjoint (that is what the case says) and cover the right
side; equal finite cardinalities force every neighbourhood to be a singleton.

\emph{The four combinations} (\texttt{SS:1534}): the proof derives
$\proj_{1,4}(\delta)=\sigma$ or $\proj_{1,3}(\delta)=\sigma$, and by the same
argument with $x_{1},x_{2}$ switched $\proj_{2,4}(\delta)=\sigma$ or
$\proj_{2,3}(\delta)=\sigma$, then says ``This completes this case''. Two of
the four combinations are the conclusion; the other two must be excluded, and
they are, because either would give $\proj_{1,2}(\delta)\subseteq\sigma$
against $\proj_{1,2}(\delta)=\cover\sigma\supsetneq\sigma$. The reverse
inclusion in the conclusion comes from stability under $\sigma$. The
``switching $x_{1}$ and $x_{2}$'' appeal needs $\cover\sigma$ symmetric.

The auxiliary relations $\zeta_{1},\zeta_{2}$ \emph{are} defined in the source,
at \texttt{SS:1486--1497}, and the proof paper reproduces those definitions
rather than reconstructing them: $\zeta_{1}$ carries eight existentially
quantified variables and four conjuncts, and the argument is a \emph{subcase
split} on whether $\zeta_{1}$ is a bridge --- not a claim that both are. An
external review was right that an earlier draft of the proof paper invoked them
without definition, and an intermediate draft supplied two short formulas that
were not the source's.

One gap at the Case~2 call: hypothesis (3) of the bridge definition for
$\xi\cap B^{4}$ requires $\cover\sigma\cap B^{2}\supsetneq\sigma\cap B^{2}$,
which is not argued. It is not needed --- if it fails, clause (d) puts both
pairs of every tuple in $\sigma$, and linkedness on a $B$ that is not a
$\sigma$-block produces $(x_{1},x_{1},x_{3},x_{3})\in\xi$ with
$(x_{1},x_{3})\notin\sigma$ directly, which is the conclusion. But the lemma
cannot be \emph{cited} on a non-bridge; the case has to be inlined. Proof
paper: inside \pkey{lem:pc-bridges-trivial}.

\subsection{\S5.5, types interaction}\label{sec:exp-55}

Both statements are correct. \texttt{LEMBridgeBetweenCongruences} omits one
direction: clause (4) is an \emph{equivalence}, and the proof establishes only
$(x_{1},x_{2})\in\sigma_{1}\Rightarrow(y_{1},y_{2})\in\sigma_{2}$. The converse
is the mirror argument, and it is where the \emph{other} half of the hypothesis
$\omega\cap\sigma_{1}=\omega\cap\sigma_{2}$ is spent. Also unstated: clause (3)
needs $\sigma_{1}\subseteq\proj_{1,2}(\delta)$ as well as the strictness
witness, and $(a,b)\notin\sigma_{2}$ follows because
$\omega\setminus\sigma_{1}=\omega\setminus\sigma_{2}$.

\texttt{LEMTwoStableIntersection} is missing one case and six steps. The
missing case is the mirror $T_{1}=\TD$, $T_{2}\in\{\TBA,\TC\}$: the hypotheses
are symmetric so it is the same argument with the indices swapped, but nothing
says so, where \texttt{SS:1805} does say ``similarly'' for the mirror of the
$\TS$ case. The six steps, all short and all discharged by
$C_{1}\cap C_{2}=\varnothing$:

\begin{enumerate}\tightlist
\item $C_{1}\cap B_{2}\subt{T_{1}}B_{1}\cap B_{2}$ needs \emph{properness}: if
  $C_{1}\cap B_{2}=B_{1}\cap B_{2}$ then $B_{1}\cap C_{2}\subseteq
  C_{1}\cap C_{2}=\varnothing$.
\item \texttt{SS:1830}: factoring returns $\subseteq$ and can collapse a proper
  containment; properness holds because $(C_{1}\cap B_{2})/\sigma_{2}=
  B_{2}/\sigma_{2}$ would make the $\sigma_{2}$-block defining $C_{2}$ meet
  $C_{1}$.
\item \texttt{SS:1817} writes the first alternative as $(B_{1}\cap B_{2})/
  \sigma_{2}=C_{2}/\sigma_{2}$ where the clause supplies only size $1$; which
  block it is follows from $B_{1}\cap C_{2}\neq\varnothing$.
\item The $\TD\times\TD$ case must apply the clause to $C_{2}$, not $B_{2}$
  (\S\ref{sec:leg-chain}).
\item The chain for $C_{2}\lll\mathbf A$ must be \emph{chosen} to extend the
  chain for $B_{2}\lll\mathbf A$ (\S\ref{sec:leg-chain}).
\item $(c_{1},c_{2})\in\omega\setminus\sigma_{1}$: membership in $\omega$
  because each $B_{i}$ lies in a single block of every dividing congruence of
  its own chain and $c_{1},c_{2}\in B_{1}\cap B_{2}$; non-membership in
  $\sigma_{1}$ because $c_{2}\in B_{1}$ and $c_{2}\,\sigma_{1}\,c_{1}\in C_{1}$
  would give $c_{2}\in C_{1}\cap C_{2}$.
\end{enumerate}

Proof paper: \pkey{lem:bridge-between-congruences}, \pkey{lem:two-stable}.

\subsection{\S5.6 and \S5.7}\label{sec:exp-5657}

Beyond $\mathrm D15$, $\mathrm D16$ and $\mathrm D17$, the remaining sixteen
proofs are sound. The steps to supply:

\begin{itemize}\tightlist
\item \texttt{LEMCenterCanBePushedIn} takes the ``full'' alternative of the
  intersection corollary without a word. ``Empty'' is excluded by
  $R\cap(B_{1}\times B_{2})\neq\varnothing$. \emph{``Size one'' needs two
  lines}: the class would have to be $[c]_{\delta}$, since $c$ has an
  $R$-partner in $B_{1}$; and then $B_{2}''=B_{2}'\cap E$ either has
  $E=[c]_{\delta}$, putting $c$ in $B_{2}''$ against the minimal choice of
  $B_{2}'$, or misses $[c]_{\delta}$ entirely, so $R\cap(B_{1}\times B_{2}'')=
  \varnothing$ against $B_{2}\subseteq B_{2}''$. Smaller: $S'=(B_{1}/\sigma)^{
  |A_{1}|}$ holds because a \emph{single} good $c$ witnesses every tuple, and
  $S''\neq\varnothing$ is where $R\cap(B_{1}\times B_{2})\neq\varnothing$ is
  spent.
\item \texttt{LEMLeftLinkedStayFull}: Case~2's restriction to $n=2^{k}$ is not
  cosmetic --- $R_{n}=(R\circ R^{-1})^{n}$ is symmetric, so
  $R_{2n}=R_{n}\circ R_{n}^{-1}$, which is what makes
  $\operatorname{LeftLinked}$ of the auxiliary relation computable in one step.
  Unstated: the auxiliary relation is subdirect over
  $(B_{1}/\sigma)\times\proj_{2}$ of itself, not over $(B_{1}/\sigma)\times
  A_{2}$; and $R_{n}\supseteq0_{A_{1}}$ by subdirectness.
\item \texttt{LEMCongruenceEitherCutOrDoNothing} leans throughout on
  $B^{2}\subseteq\omega$ --- every $B_{i}$ in the chain lies inside a single
  block of its own dividing congruence --- which is never stated and is used at
  least four times. Also unstated: $|B/\sigma|>1$, and that a path witnessing
  linkedness must contain a single step crossing $\sigma$-classes.
\item \texttt{LEMFactorByDelta}: Case~2 is a chain of six steps of which the
  source states three, and the \textbf{``moreover'' clause is never argued} ---
  it holds because Case~2 gives $\delta\cap B^{2}\subseteq\delta'\cap B^{2}=
  \sigma\cap B^{2}$ and Case~1 never produces the type-$T$ alternative. That
  clause is the saturation identity, and it is what \S\ref{sec:D17} turns on.
\item \texttt{LEMPropagation}: clause (ft) is discharged by citing the
  factorisation lemma, which is stated only for $T\in\{\TPC,\TL\}$; the
  $\{\TBA,\TC,\TS\}$ cases need clause (fs) and a sentence. And clause (fm)'s
  $t=1$ base case is where the $\TS$-freeness hypothesis is \emph{spent} --- it
  kills the factorisation lemma's middle alternative, which a multi-type cannot
  absorb. Nothing says so.
\item \texttt{CORPropagateToRelations}(e) needs two unstated steps:
  $C\circ\delta=B\circ\delta$ is excluded because $\delta\subseteq\sigma$ puts
  $C\circ\delta$ inside one $\sigma$-block while $|B/\sigma|>1$; and the
  $\subt{\TS}$ alternative is excluded because the witness also lies in one
  $\sigma$-block.
\item \texttt{LEMMultiplyByAllLinear} invokes the stable-intersection theorem
  (\texttt{SS:3057}) for its Case~1/Case~2 dichotomy. An earlier reading of
  this entry said no theorem is consumed, on the ground that walking down the
  chain and asking where the intersection with $B_{2}$ becomes empty is
  elementary. \textbf{That is right about the dichotomy and wrong about what
  Case~2 asserts.} Case~2 does not merely produce a step; it produces a step
  \emph{of type $\TD$}, and the inductive hypothesis is then applied to it,
  which needs it to be of multi-type $T$. Two things are therefore owed and
  neither is given: that the step is not of type $\TBA$, $\TC$ or $\TS$ ---
  provable for $\TS$ from \pkey{lem:intersection-good}(s), but not for the
  other two, since $\subt{\TBA,\TC}$ there is a conjunction and one absorbing
  step supplies half of it --- and, when $T\in\{\TL,\TPC\}$, that the
  irreducible congruence carried by a $\TD$ step is of type $T$ rather than the
  other one. The second is a missing hypothesis in the induction rather than a
  suppressed step. The stable-intersection appeal is presumably aimed at the
  first, but the theorem's conclusion is about a critically inconsistent family
  having equal types, and the family here is not put in that form.
  Proof paper: \pkey{haz:multiply-citation}, \pkey{cl:multiply-step}.
\item \texttt{LEMPCOnTheTopIsEasy} rests on two things it does not state: that
  an idempotent algebra is polynomially complete iff every \emph{reflexive}
  subalgebra of a power is a conjunction of equalities, and that
  $\mathbf A/\sigma$ is BA and centre free, without which the step
  $\{a_{i}\}\subt{\TPC}\mathbf A/\sigma$ is not licensed. Proof paper:
  \pkey{haz:on-top}.
\item Typographical: \texttt{LEMUbiquity}'s ``block $D$ of $B/\delta$'' should
  read ``block $D$ of $\delta$''; \texttt{LEMCenter\brk CanBePushedIn}'s $\neq$ at
  \texttt{SS:2014} should be $\not\subseteq$;
  \texttt{CORParallelogramPropertyForD}'s undefined $C$ at \texttt{SS:1000}
  should be $C_{1}$; \texttt{LEMMaximalMultExtention}'s
  $\cover{\delta_{i}}\supseteq B_{1}^{2}$ at \texttt{SS:3128} should be
  $\cover{\delta_{i}}\supseteq(B_{1}/\sigma)^{2}$; and the instance-fragmented
  definition writes $X_{1}$ twice where the second should be $X_{2}$.
\end{itemize}

\subsection{\texttt{LEMPreserveLinkdness}: a flag withdrawn}\label{sec:exp-preserve}

Flagged in an earlier survey on the grounds that the workhorse
\texttt{LEMPreserve\brk Linkdness\brk OneStepAUX} assumes both
$R\cap(B_{1}\times C_{2})\neq\varnothing$ and $R\cap(C_{1}\times B_{2})\neq
\varnothing$ while the target lemma assumes only $R\cap(B_{1}\times B_{2})\neq
\varnothing$. \textbf{Not a defect}: the flag conflates the auxiliary lemma
with \texttt{LEMPreserveLinkdnessOneStep}, which is what the four-line proof
actually invokes and which needs only $R\cap(B_{1}\times B_{2})\neq\varnothing$
and $S\cap(C_{1}\times B_{2})\neq\varnothing$.

The four-line proof is correct and unpacks to a two-phase argument the proof
paper should write out. Decompose both multi-type reductions into chains.
\emph{Phase 1}: shrink coordinate 1 along its chain with $B_{2}$ fixed; at each
step the needed $R\cap(B_{1}^{(j)}\times B_{2})\neq\varnothing$ is the previous
step's conclusion, and the needed $S\cap(B_{1}^{(j+1)}\times B_{2})\neq
\varnothing$ follows from $S\cap(C_{1}\times C_{2})\neq\varnothing$ by
monotonicity. This yields $R\cap(C_{1}\times B_{2})\neq\varnothing$.
\emph{Phase 2}: the same on the transpose $R^{-1}$, whose rectangular closure
is $S^{-1}$, shrinking coordinate 2 with $C_{1}$ fixed --- legitimate because
$C_{1}\lll A_{1}$. Proof paper: \pkey{lem:preserve-linked}.

The withdrawal is of this flag only. A later reading of the auxiliary
lemma's \emph{own} proof found a real defect in its Case~2 --- the
cross-component step is asserted, not shown --- recorded as
$\mathrm D21$ (\S\ref{sec:D21}).
